\documentclass[11pt,a4paper]{article}

% ============================================================
% PACOTES
% ============================================================

\usepackage[utf8]{inputenc}
\usepackage[T1]{fontenc}
\usepackage[portuguese]{babel}
\usepackage{lmodern}

\usepackage[
a4paper,
top=1.8cm,
bottom=1.8cm,
left=1.8cm,
right=1.8cm
]{geometry}

\usepackage{amsmath,amssymb}
\usepackage{array}
\usepackage{booktabs}
\usepackage{tabularx}
\usepackage{multirow}
\usepackage{multicol}
\usepackage{enumitem}
\usepackage{float}
\usepackage{xcolor}
\usepackage{graphicx}

% ------------------------------------------------------------
% SIUNITX
% ------------------------------------------------------------
\usepackage{siunitx}
\sisetup{output-decimal-marker={,}}

% Fórmulas químicas
\usepackage[version=4]{mhchem}

% TikZ
\usepackage{tikz}
\usetikzlibrary{
	arrows.meta,
	positioning,
	calc,
	shapes.geometric
}

\usepackage{pgfplots}
\pgfplotsset{compat=1.18}

% Caixas
\usepackage[most]{tcolorbox}

% Cabeçalhos
\usepackage{fancyhdr}
\usepackage{lastpage}

\usepackage{hyperref}

% ============================================================
% CORES
% ============================================================

\definecolor{eublue}{HTML}{145A86}
\definecolor{eudark}{HTML}{17324D}
\definecolor{eulightblue}{HTML}{E4F2FA}
\definecolor{eugreen}{HTML}{3F7D58}
\definecolor{eulightgreen}{HTML}{E8F4EC}
\definecolor{euorange}{HTML}{E28A2B}
\definecolor{eulightorange}{HTML}{FFF0DE}
\definecolor{eured}{HTML}{B33A3A}
\definecolor{eupurple}{HTML}{67518A}
\definecolor{eugray}{HTML}{F2F4F6}

% ============================================================
% CONFIGURAÇÕES
% ============================================================

\setlength{\parindent}{0pt}
\setlength{\parskip}{5pt}

\renewcommand{\arraystretch}{1.35}

\hypersetup{
	colorlinks=true,
	linkcolor=eublue,
	urlcolor=eublue
}

% ============================================================
% CABEÇALHO E RODAPÉ
% ============================================================

\pagestyle{fancy}
\fancyhf{}

\fancyhead[L]{
	\textbf{\textcolor{eublue}{EUFISICA exercícios}}
}

\fancyhead[R]{
	\textcolor{eudark}{Resolução -- Física e Química A -- 10.º ano}
}

\fancyfoot[L]{Prof. José Gonçalves}
\fancyfoot[C]{EUFISICA}
\fancyfoot[R]{Página \thepage\ de \pageref{LastPage}}

\renewcommand{\headrulewidth}{0.6pt}
\renewcommand{\footrulewidth}{0.4pt}

% ============================================================
% CAIXAS
% ============================================================

\newcounter{exercicio}

\newenvironment{resolucao}[1]
{
	\refstepcounter{exercicio}
	\begin{tcolorbox}[
		enhanced,
		breakable,
		colback=white,
		colframe=eublue,
		boxrule=0.8pt,
		arc=2mm,
		left=3mm,
		right=3mm,
		top=2mm,
		bottom=2mm,
		title={
			\textbf{Exercício \theexercicio}
			\hspace{1mm}--\hspace{1mm}#1
		},
		coltitle=white,
		colbacktitle=eublue,
		fonttitle=\bfseries
		]
	}
	{
	\end{tcolorbox}
}

\newtcolorbox{resultado}{
	enhanced,
	colback=eulightgreen,
	colframe=eugreen,
	fontupper=\bfseries,
	arc=2mm
}

\newtcolorbox{nota}{
	enhanced,
	breakable,
	colback=eulightorange,
	colframe=euorange,
	title=\textbf{Nota},
	fonttitle=\bfseries,
	arc=2mm
}

\newtcolorbox{recorda}{
	enhanced,
	breakable,
	colback=eulightblue,
	colframe=eublue,
	title=\textbf{Expressões úteis},
	fonttitle=\bfseries,
	arc=2mm
}

% ============================================================
% DOCUMENTO
% ============================================================

\begin{document}
	
	% ============================================================
	% TÍTULO
	% ============================================================
	
	\begin{tcolorbox}[
		enhanced,
		colback=eudark,
		colframe=eudark,
		arc=3mm,
		boxrule=0pt,
		left=5mm,
		right=5mm,
		top=5mm,
		bottom=5mm
		]
		
		\begin{center}
			
			{\Huge\bfseries\color{white} EUFISICA exercícios}
			
			\vspace{2mm}
			
			{\Large\bfseries\color{white}
				Resolução da Ficha de Trabalho}
			
			\vspace{2mm}
			
			{\large\color{eulightblue}
				Quantidade de matéria • Espetros atómicos • Átomo de hidrogénio}
			
			{\large\color{eulightblue}
				Energia dos eletrões • Espetro fotoeletrónico • Distribuição eletrónica}
			
			\vspace{4mm}
			
			{\large\color{white}\textbf{Física e Química A -- 10.º ano}}
			
			\vspace{2mm}
			
			{\color{white}\textbf{Prof. José Gonçalves}}
			
		\end{center}
		
	\end{tcolorbox}
	
	% ============================================================
	% CONSTANTES
	% ============================================================
	
	\begin{recorda}
		
		\[
		N_A=\SI{6.022e23}{\per\mole}
		\]
		
		\[
		h=\SI{6.626e-34}{\joule\second}
		\]
		
		\[
		c=\SI{3.00e8}{\meter\per\second}
		\]
		
		\[
		1\,\mathrm{eV}
		=
		\SI{1.602e-19}{J}
		\]
		
		\[
		N=nN_A
		\]
		
		\[
		n=\frac{m}{M}
		\]
		
		\[
		E_{\text{fotão}}=h f
		\]
		
		\[
		c=\lambda f
		\]
		
		\[
		E_{\text{fotão}}
		=
		\frac{hc}{\lambda}
		\]
		
		Para o átomo de hidrogénio:
		
		\[
		\boxed{
			E_n=-\frac{2,18\times10^{-18}}{n^2}\ \mathrm{J}
		}
		\]
		
		ou
		
		\[
		\boxed{
			E_n=-\frac{13,6}{n^2}\ \mathrm{eV}
		}
		\]
		
	\end{recorda}
	
	% ============================================================
	% 1
	% ============================================================
	
	\section*{\textcolor{eublue}{1. Quantidade de matéria e massa molar}}
	
	\begin{resolucao}{Quantidade de matéria}
		
		Utiliza-se:
		
		\[
		N=nN_A
		\]
		
		Assim:
		
		\[
		N(1,00\,\mathrm{mol})
		=
		1,00\times6,022\times10^{23}
		\]
		
		\[
		N=6,022\times10^{23}
		\]
		
		Para \(\SI{0.500}{mol}\):
		
		\[
		N=
		0,500\times6,022\times10^{23}
		\]
		
		\[
		N=3,011\times10^{23}
		\]
		
		Para \(\SI{2.00}{mol}\):
		
		\[
		N=
		2,00\times6,022\times10^{23}
		\]
		
		\[
		N=1,2044\times10^{24}
		\]
		
		Nos restantes casos:
		
		\[
		n=\frac{N}{N_A}
		\]
		
		\[
		n=
		\frac{3,011\times10^{23}}
		{6,022\times10^{23}}
		=
		0,500\,\mathrm{mol}
		\]
		
		e
		
		\[
		n=
		\frac{1,2044\times10^{24}}
		{6,022\times10^{23}}
		=
		2,00\,\mathrm{mol}
		\]
		
		\begin{center}
			
			\begin{tabular}{ccc}
				\toprule
				\textbf{Quantidade de matéria}
				&
				\textbf{Número de entidades}
				&
				\textbf{Entidade}
				\\
				\midrule
				\(\SI{1.00}{mol}\)
				&
				\(6,022\times10^{23}\)
				&
				átomos
				\\
				
				\(\SI{0.500}{mol}\)
				&
				\(3,011\times10^{23}\)
				&
				moléculas
				\\
				
				\(\SI{2.00}{mol}\)
				&
				\(1,2044\times10^{24}\)
				&
				iões
				\\
				
				\(\SI{0.500}{mol}\)
				&
				\(3,011\times10^{23}\)
				&
				moléculas
				\\
				
				\(\SI{2.00}{mol}\)
				&
				\(1,2044\times10^{24}\)
				&
				átomos
				\\
				\bottomrule
			\end{tabular}
			
		\end{center}
		
	\end{resolucao}
	
	% ============================================================
	
	\begin{resolucao}{Massa molar}
		
		\textbf{a) Água, \(\ce{H2O}\)}
		
		\[
		M(\ce{H2O})
		=
		2M(\ce{H})+M(\ce{O})
		\]
		
		\[
		M(\ce{H2O})
		=
		2(1,01)+16,00
		\]
		
		\[
		\boxed{
			M(\ce{H2O})=\SI{18.02}{g\per mol}
		}
		\]
		
		\vspace{2mm}
		
		\textbf{b) Dióxido de carbono, \(\ce{CO2}\)}
		
		\[
		M(\ce{CO2})
		=
		12,01+2(16,00)
		\]
		
		\[
		\boxed{
			M(\ce{CO2})=\SI{44.01}{g\per mol}
		}
		\]
		
		\vspace{2mm}
		
		\textbf{c) Amoníaco, \(\ce{NH3}\)}
		
		\[
		M(\ce{NH3})
		=
		14,01+3(1,01)
		\]
		
		\[
		\boxed{
			M(\ce{NH3})=\SI{17.03}{g\per mol}
		}
		\]
		
		\vspace{2mm}
		
		\textbf{d) Cloreto de sódio, \(\ce{NaCl}\)}
		
		\[
		M(\ce{NaCl})
		=
		22,99+35,45
		\]
		
		\[
		\boxed{
			M(\ce{NaCl})=\SI{58.44}{g\per mol}
		}
		\]
		
		\vspace{2mm}
		
		\textbf{e) Carbonato de cálcio, \(\ce{CaCO3}\)}
		
		\[
		M(\ce{CaCO3})
		=
		40,08+12,01+3(16,00)
		\]
		
		\[
		\boxed{
			M(\ce{CaCO3})=\SI{100.09}{g\per mol}
		}
		\]
		
		\vspace{2mm}
		
		\textbf{f) Sulfato de alumínio, \(\ce{Al2(SO4)3}\)}
		
		\[
		M=
		2(26,98)
		+
		3[32,06+4(16,00)]
		\]
		
		\[
		M=
		53,96+3(96,06)
		\]
		
		\[
		\boxed{
			M(\ce{Al2(SO4)3})
			=
			\SI{342.14}{g\per mol}
		}
		\]
		
	\end{resolucao}
	
	% ============================================================
	
	\begin{resolucao}{Da massa ao número de moléculas}
		
		Para a água:
		
		\[
		M(\ce{H2O})=\SI{18.02}{g\per mol}
		\]
		
		\textbf{a)}
		
		\[
		\boxed{
			M(\ce{H2O})=\SI{18.02}{g\per mol}
		}
		\]
		
		\textbf{b)}
		
		\[
		n=\frac{m}{M}
		\]
		
		\[
		n=
		\frac{18,0}{18,02}
		\]
		
		\[
		\boxed{
			n\approx\SI{0.999}{mol}
		}
		\]
		
		ou, aproximadamente,
		
		\[
		n\approx\SI{1.00}{mol}.
		\]
		
		\textbf{c)}
		
		\[
		N=nN_A
		\]
		
		\[
		N=
		0,999\times6,022\times10^{23}
		\]
		
		\[
		\boxed{
			N\approx6,02\times10^{23}
			\text{ moléculas}
		}
		\]
		
		\textbf{d)}
		
		Cada molécula de água possui dois átomos de hidrogénio.
		
		Assim:
		
		\[
		N(\ce{H})
		=
		2N(\ce{H2O})
		\]
		
		\[
		N(\ce{H})
		=
		2(6,02\times10^{23})
		\]
		
		\[
		\boxed{
			N(\ce{H})
			\approx1,20\times10^{24}
			\text{ átomos}
		}
		\]
		
	\end{resolucao}
	
	% ============================================================
	
	\begin{resolucao}{Um composto do quotidiano}
		
		A cafeína apresenta fórmula:
		
		\[
		\ce{C8H10N4O2}
		\]
		
		\textbf{a) Massa molar}
		
		\[
		M=
		8(12,01)+10(1,01)+4(14,01)+2(16,00)
		\]
		
		\[
		M=
		96,08+10,10+56,04+32,00
		\]
		
		\[
		\boxed{
			M=\SI{194.22}{g\per mol}
		}
		\]
		
		\textbf{b) Quantidade de matéria}
		
		Primeiro:
		
		\[
		\SI{80}{mg}
		=
		\SI{0.080}{g}
		\]
		
		Logo,
		
		\[
		n=
		\frac{0,080}{194,22}
		\]
		
		\[
		\boxed{
			n\approx4,12\times10^{-4}\,\mathrm{mol}
		}
		\]
		
		\textbf{c) Número de moléculas}
		
		\[
		N=nN_A
		\]
		
		\[
		N=
		4,12\times10^{-4}
		\times6,022\times10^{23}
		\]
		
		\[
		\boxed{
			N\approx2,48\times10^{20}
			\text{ moléculas}
		}
		\]
		
	\end{resolucao}
	
	% ============================================================
	% 2
	% ============================================================
	
	\newpage
	
	\section*{\textcolor{eublue}{2. Espetros de emissão e de absorção}}
	
	\begin{resolucao}{Emissão ou absorção?}
		
		\begin{enumerate}[label=\alph*)]
			
			\item
			
			\[
			E_4\rightarrow E_2
			\]
			
			O eletrão passa de um nível de maior energia para um nível de menor
			energia.
			
			\[
			\boxed{\text{Emissão}}
			\]
			
			\item
			
			\[
			E_1\rightarrow E_3
			\]
			
			O eletrão passa para um nível de maior energia.
			
			\[
			\boxed{\text{Absorção}}
			\]
			
			\item Uma risca luminosa sobre fundo escuro é característica de um:
			
			\[
			\boxed{\text{espetro de emissão}}
			\]
			
			\item Uma risca escura num espetro contínuo corresponde a:
			
			\[
			\boxed{\text{absorção}}
			\]
			
			\item O átomo recebe um fotão cuja energia corresponde à diferença
			entre dois níveis, permitindo ao eletrão transitar para um nível superior.
			
			\[
			\boxed{\text{Absorção}}
			\]
			
		\end{enumerate}
		
	\end{resolucao}
	
	% ============================================================
	
	\begin{resolucao}{Identificação de elementos por espetroscopia}
		
		\textbf{a)}
		
		A amostra X apresenta linhas coincidentes com as linhas do hidrogénio:
		
		\[
		410,\ 434,\ 486,\ 656\ \mathrm{nm}
		\]
		
		e com linhas do sódio próximas de:
		
		\[
		568\ \mathrm{nm}
		\qquad\text{e}\qquad
		589\ \mathrm{nm}.
		\]
		
		Logo:
		
		\[
		\boxed{
			\text{A amostra contém hidrogénio e sódio.}
		}
		\]
		
		\textbf{b)}
		
		Cada elemento possui um conjunto próprio de níveis eletrónicos de
		energia. As diferenças entre esses níveis também são características
		do elemento.
		
		Como:
		
		\[
		E_{\text{fotão}}
		=
		|\Delta E|
		\]
		
		cada elemento origina um conjunto característico de comprimentos de
		onda.
		
		Por isso, o espetro funciona como uma espécie de
		\textbf{impressão digital do elemento}.
		
		\textbf{c)}
		
		Não se espera que dois elementos diferentes possuam espetros
		completamente iguais.
		
		Podem existir algumas linhas com comprimentos de onda muito próximos,
		mas o conjunto completo de linhas é característico de cada elemento.
		
	\end{resolucao}
	
	% ============================================================
	% 3
	% ============================================================
	
	\section*{\textcolor{eublue}{3. Energia dos fotões}}
	
	\begin{resolucao}{Energia de um fotão}
		
		Dado:
		
		\[
		\lambda=\SI{500}{nm}
		\]
		
		\textbf{a) Conversão para metros}
		
		\[
		\SI{1}{\nano\meter} = \SI{1e-9}{\meter}
		\]
		
		logo:
		
		\[
		\SI{500}{\nano\meter}
		=
		\SI{500e-9}{\meter}
		\]
		
		\[
		\boxed{
			\lambda=\SI{5.00e-7}{m}
		}
		\]
		
		\textbf{b) Frequência}
		
		\[
		c=\lambda f
		\]
		
		\[
		f=\frac{c}{\lambda}
		\]
		
		\[
		\nu=
		\frac{3,00\times10^8}
		{5,00\times10^{-7}}
		\]
		
		\[
		\boxed{
			\nu=\SI{6.00e14}{Hz}
		}
		\]
		
		\textbf{c) Energia de um fotão}
		
		\[
		E=h f
		\]
		
		\[
		E=
		6,626\times10^{-34}
		\times6,00\times10^{14}
		\]
		
		\[
		\boxed{
			E\approx\SI{3.98e-19}{J}
		}
		\]
		
		\textbf{d) Energia correspondente a um mol de fotões}
		
		\[
		E_{\mathrm{mol}}
		=
		EN_A
		\]
		
		\[
		E_{\mathrm{mol}}
		=
		3,98\times10^{-19}
		\times
		6,022\times10^{23}
		\]
		
		\[
		E_{\mathrm{mol}}
		\approx
		\SI{2,39e5}{\joule\per\mole}
		\]
		
		ou:
		
		\[
		\boxed{
			E_{\mathrm{mol}}
			\approx
			\SI{239}{kJ\per mol}
		}
		\]
		
	\end{resolucao}
	
	% ============================================================
	
	\begin{resolucao}{Comparar fotões}
		
		Como:
		
		\[
		f=\frac{c}{\lambda}
		\]
		
		quanto menor for \(\lambda\), maior será \(f\).
		
		Como:
		
		\[
		E=hf
		\]
		
		quanto maior for a frequência, \(f\), maior será a energia.
		
		Os comprimentos de onda são:
		
		\[
		\lambda_A=\SI{400}{nm}
		\]
		
		\[
		\lambda_B=\SI{550}{nm}
		\]
		
		\[
		\lambda_C=\SI{650}{nm}
		\]
		
		\textbf{a)}
		
		O menor comprimento de onda é o de A.
		
		\[
		\boxed{\text{A possui maior frequência.}}
		\]
		
		\textbf{b)}
		
		C apresenta o maior comprimento de onda.
		
		\[
		\boxed{\text{C possui menor energia.}}
		\]
		
		\textbf{c)}
		
		Por ordem crescente de energia:
		
		\[
		\boxed{
			C<B<A
		}
		\]
		
		\textbf{d)}
		
		Pode justificar-se diretamente por:
		
		\[
		\boxed{
			E=\frac{hc}{\lambda}
		}
		\]
		
		Logo, a energia é inversamente proporcional ao comprimento de onda.
		
	\end{resolucao}
	
	% ============================================================
	% 4 HIDROGÉNIO
	% ============================================================
	
	\newpage
	
	\section*{\textcolor{eublue}{4. Espetro e níveis de energia do átomo de hidrogénio}}
	
	\begin{resolucao}{Energia dos níveis do hidrogénio}
		
		Utiliza-se:
		
		\[
		E_n=
		-\frac{2,18\times10^{-18}}{n^2}\ \mathrm{J}
		\]
		
		\textbf{a) Para \(n=1\)}
		
		\[
		E_1=
		-\frac{2,18\times10^{-18}}{1^2}
		\]
		
		\[
		\boxed{
			E_1=\SI{-2,18e-18}{\joule}
		}
		\]
		
		\textbf{b) Para \(n=2\)}
		
		\[
		E_2=
		-\frac{2,18\times10^{-18}}{2^2}
		\]
		
		\[
		E_2=
		-\frac{2,18\times10^{-18}}{4}
		\]
		
		\[
		\boxed{
			E_2=\SI{-5,45e-19}{\joule}
		}
		\]
		
		\textbf{c) Para \(n=3\)}
		
		\[
		E_3=
		-\frac{2,18\times10^{-18}}{9}
		\]
		
		\[
		\boxed{
			E_3\approx\SI{-2,42e-19}{\joule}
		}
		\]
		
		\textbf{d) Para \(n=4\)}
		
		\[
		E_4=
		-\frac{2,18\times10^{-18}}{16}
		\]
		
		\[
		\boxed{
			E_4\approx\SI{-1,36e-19}{\joule}
		}
		\]
		
		Verifica-se que, à medida que \(n\) aumenta, a energia aproxima-se de
		zero.
		
	\end{resolucao}
	
	% ============================================================
	
	\begin{resolucao}{Transição eletrónica \(n=4\rightarrow n=2\)}
		
		\textbf{a)}
		
		\[
		E_4=
		-\frac{2,18\times10^{-18}}{16}
		\]
		
		\[
		\boxed{
			E_4=\SI{-1,36e-19}{\joule}
		}
		\]
		
		\textbf{b)}
		
		\[
		E_2=
		-\frac{2,18\times10^{-18}}{4}
		\]
		
		\[
		\boxed{
			E_2=\SI{-5,45e-19}{\joule}
		}
		\]
		
		\textbf{c) Variação de energia}
		
		\[
		\Delta E=E_f-E_i
		\]
		
		\[
		\Delta E
		=
		E_2-E_4
		\]
		
		\[
		\Delta E
		=
		-5,45\times10^{-19}
		-
		(-1,36\times10^{-19})
		\]
		
		\[
		\boxed{
			\Delta E\approx\SI{-4,09e-19}{\joule}
		}
		\]
		
		\textbf{d)}
		
		Como:
		
		\[
		\Delta E<0
		\]
		
		o átomo perdeu energia.
		
		\[
		\boxed{\text{Ocorre emissão de um fotão.}}
		\]
		
		\textbf{e)}
		
		\[
		E_{\text{fotão}}=|\Delta E|
		\]
		
		\[
		\boxed{
			E_{\text{fotão}}
			=
			\SI{4,09e-19}{\joule}
		}
		\]
		
		\textbf{f) Frequência}
		
		\[
		E=hf
		\]
		
		\[
		f=\frac{E}{h}
		\]
		
		\[
		f=
		\frac{4,09\times10^{-19}}
		{6,626\times10^{-34}}
		\]
		
		\[
		\boxed{
			f\approx\SI{6,17e14}{\hertz}
		}
		\]
		
		\textbf{g) Comprimento de onda}
		
		\[
		\lambda=\frac{c}{\nu}
		\]
		
		\[
		\lambda=
		\frac{3,00\times10^8}
		{6,17\times10^{14}}
		\]
		
		\[
		\lambda\approx\SI{4,86e-7}{\meter}
		\]
		
		\[
		\boxed{
			\lambda\approx\SI{486}{nm}
		}
		\]
		
		\textbf{h)}
		
		Este comprimento de onda pertence à região:
		
		\[
		\boxed{\text{visível, na região azul--verde}}
		\]
		
		Esta linha corresponde à linha \(H_\beta\) da série de Balmer.
		
	\end{resolucao}
	
	% ============================================================
	
	\begin{resolucao}{Outra risca do hidrogénio}
		
		Transição:
		
		\[
		n=3\rightarrow n=2
		\]
		
		\[
		E_3=
		\SI{-2,42e-19}{J}
		\]
		
		\[
		E_2=
		\SI{-5,45e-19}{J}
		\]
		
		\textbf{a)}
		
		A transição deve ser representada por uma seta descendente de
		\(n=3\) para \(n=2\).
		
		\textbf{b)}
		
		\[
		\Delta E
		=
		E_2-E_3
		\]
		
		\[
		\Delta E=
		-5,45\times10^{-19}
		-
		(-2,42\times10^{-19})
		\]
		
		\[
		\boxed{
			\Delta E\approx\SI{-3,03e-19}{\joule}
		}
		\]
		
		\textbf{c)}
		
		\[
		E_{\text{fotão}}
		=
		|\Delta E|
		\]
		
		\[
		\boxed{
			E_{\text{fotão}}
			=
			\SI{3,03e-19}{\joule}
		}
		\]
		
		\textbf{d)}
		
		\[
		\lambda
		=
		\frac{hc}{E}
		\]
		
		\[
		\lambda
		=
		\frac{
			6,626\times10^{-34}
			\times3,00\times10^8
		}{
			3,03\times10^{-19}
		}
		\]
		
		\[
		\lambda
		\approx\SI{6,57e-7}{\meter}
		\]
		
		\[
		\boxed{
			\lambda\approx\SI{657}{nm}
		}
		\]
		
		Com valores mais precisos para os níveis obtém-se aproximadamente:
		
		\[
		\boxed{
			\lambda\approx\SI{656}{nm}
		}
		\]
		
		\textbf{e)}
		
		Este valor coincide com a risca vermelha do espetro do hidrogénio,
		próxima de:
		
		\[
		\boxed{
			\SI{656.3}{nm}
		}
		\]
		
	\end{resolucao}
	
	% ============================================================
	
	\begin{resolucao}{Absorção de radiação}
		
		Transição:
		
		\[
		n=1\rightarrow n=3
		\]
		
		\textbf{a)}
		
		\[
		\boxed{
			E_1=\SI{-2,18e-18}{\joule}
		}
		\]
		
		\textbf{b)}
		
		\[
		E_3=
		-\frac{2,18\times10^{-18}}{9}
		\]
		
		\[
		\boxed{
			E_3=\SI{-2,42e-19}{\joule}
		}
		\]
		
		\textbf{c)}
		
		\[
		\Delta E
		=
		E_3-E_1
		\]
		
		\[
		\Delta E
		=
		-2,42\times10^{-19}
		-
		(-2,18\times10^{-18})
		\]
		
		\[
		\boxed{
			\Delta E
			\approx
			\SI{1,94e-18}{\joule}
		}
		\]
		
		Portanto:
		
		\[
		\boxed{
			E_{\text{fotão,min}}
			\approx
			\SI{1,94e-18}{\joule}
		}
		\]
		
		\textbf{d)}
		
		\[
		f=\frac{E}{h}
		\]
		
		\[
		f=
		\frac{1,94\times10^{-18}}
		{6,626\times10^{-34}}
		\]
		
		\[
		\boxed{
			f\approx\SI{2,92e15}{\hertz}
		}
		\]
		
		\textbf{e)}
		
		Os níveis de energia no átomo estão quantizados.
		
		Para ocorrer a transição \(n=1\rightarrow n=3\), o fotão deve fornecer
		uma energia igual à diferença:
		
		\[
		E_3-E_1.
		\]
		
		Um fotão de energia inferior não fornece a energia necessária para essa
		transição.
		
	\end{resolucao}
	
	% ============================================================
	% DESAFIO HIDROGÉNIO
	% ============================================================
	
	\begin{tcolorbox}[
		enhanced,
		breakable,
		colback=eulightorange,
		colframe=euorange,
		title=\textbf{Resolução do desafio -- Ionização a partir de \(n=2\)},
		fonttitle=\bfseries
		]
		
		O eletrão encontra-se inicialmente em:
		
		\[
		n=2
		\]
		
		com energia:
		
		\[
		E_2=\SI{-5,45e-19}{\joule}
		\]
		
		Um eletrão completamente separado do átomo possui:
		
		\[
		E_\infty=0.
		\]
		
		Assim:
		
		\[
		\Delta E
		=
		E_\infty-E_2
		\]
		
		\[
		\Delta E
		=
		0-(-5,45\times10^{-19})
		\]
		
		\[
		\boxed{
			E_{\text{ionização}}
			=
			\SI{5,45e-19}{\joule}
		}
		\]
		
		\textbf{a) Em joule por eletrão}
		
		\[
		\boxed{
			\SI{5,45e-19}{\joule}
		}
		\]
		
		\textbf{b) Em eletrão-volt}
		
		\[
		E=
		\frac{5,45\times10^{-19}}
		{1,602\times10^{-19}}
		\]
		
		\[
		\boxed{
			E\approx\SI{3,40}{\electronvolt}
		}
		\]
		
		\textbf{c) Em \(\si{\kilo\joule\per mol}\)}
		
		\[
		E_{\mathrm{mol}}
		=
		5,45\times10^{-19}
		\times6,022\times10^{23}
		\]
		
		\[
		E_{\mathrm{mol}}
		=
		\SI{3,28e5}{\joule\per\mole}
		\]
		
		\[
		\boxed{
			E_{\mathrm{mol}}
			\approx
			\SI{328}{kJ\per mol}
		}
		\]
		
	\end{tcolorbox}
	
	% ============================================================
	% FOTOELETRÓNICO
	% ============================================================
	
	\newpage
	
	\section*{\textcolor{eublue}{5. Energias de remoção e espetro fotoeletrónico}}
	
	\begin{resolucao}{Energia de remoção}
		
		Tem-se:
		
		\[
		E_{\text{fotão}}
		=
		E_{\text{remoção}}
		+
		E_c
		\]
		
		Portanto:
		
		\[
		E_{\text{remoção}}
		=
		E_{\text{fotão}}-E_c
		\]
		
		\[
		E_{\text{remoção}}
		=
		2,40\times10^{-17}
		-
		7,0\times10^{-18}
		\]
		
		\[
		\boxed{
			E_{\text{remoção}}
			=
			\SI{1,70e-17}{\joule}
		}
		\]
		
		\textbf{b) Conversão para eV}
		
		\[
		E=
		\frac{1,70\times10^{-17}}
		{1,602\times10^{-19}}
		\]
		
		\[
		\boxed{
			E_{\text{remoção}}
			\approx\SI{106}{\electronvolt}
		}
		\]
		
		\textbf{c) Energia por mole}
		
		\[
		E_{\mathrm{mol}}
		=
		1,70\times10^{-17}
		\times
		6,022\times10^{23}
		\]
		
		\[
		E_{\mathrm{mol}}
		=
		\SI{1,024e7}{\joule\per\mole}
		\]
		
		\[
		\boxed{
			E_{\mathrm{mol}}
			\approx
			\SI{1,02e4}{\kilo\joule\per\mole}
		}
		\]
		
	\end{resolucao}
	
	% ============================================================
	
	\begin{resolucao}{Interpretar um espetro fotoeletrónico}
		
		A distribuição eletrónica indicada é:
		
		\[
		1s^2\,2s^2\,2p^6\,3s^1
		\]
		
		\textbf{a)}
		
		Existem quatro subníveis ocupados:
		
		\[
		\boxed{
			1s,\quad2s,\quad2p,\quad3s
		}
		\]
		
		\textbf{b)}
		
		Os eletrões \(1s\) encontram-se mais próximos do núcleo e estão mais
		fortemente ligados.
		
		\[
		\boxed{1s}
		\]
		
		é o subnível de maior energia de remoção.
		
		\textbf{c)}
		
		Os eletrões \(3s\) são os mais externos.
		
		\[
		\boxed{3s}
		\]
		
		apresenta os eletrões mais facilmente removíveis.
		
		\textbf{d)}
		
		O subnível \(2p\) possui:
		
		\[
		6\text{ eletrões}
		\]
		
		enquanto \(2s\) possui:
		
		\[
		2\text{ eletrões}.
		\]
		
		Por isso, a intensidade relativa do pico \(2p\) é aproximadamente três
		vezes maior.
		
		\textbf{e)}
		
		Número total de eletrões:
		
		\[
		2+2+6+1=11
		\]
		
		\[
		\boxed{11\text{ eletrões}}
		\]
		
		\textbf{f)}
		
		Um átomo neutro com 11 eletrões apresenta:
		
		\[
		Z=11
		\]
		
		e corresponde ao:
		
		\[
		\boxed{\ce{Na}}
		\]
		
		sódio.
		
	\end{resolucao}
	
	% ============================================================
	
	\begin{resolucao}{Energias de ionização sucessivas}
		
		Os valores são:
		
		\[
		738,\quad1451,\quad7733,\quad10543,\quad13630
		\]
		
		em \(\si{\kilo\joule\per mol}\).
		
		\textbf{a)}
		
		O maior aumento ocorre entre a segunda e a terceira energias:
		
		\[
		1451\rightarrow7733
		\]
		
		\[
		\boxed{\text{entre a remoção do 2.º e do 3.º eletrão}}
		\]
		
		\textbf{b)}
		
		Os dois primeiros eletrões são relativamente mais fáceis de remover.
		
		O terceiro eletrão já pertence a uma camada interna.
		
		Logo:
		
		\[
		\boxed{
			\text{o magnésio possui dois eletrões de valência}
		}
		\]
		
		\textbf{c)}
		
		Como:
		
		\[
		Z(\ce{Mg})=12
		\]
		
		a distribuição eletrónica é:
		
		\[
		\boxed{
			1s^2\,2s^2\,2p^6\,3s^2
		}
		\]
		
		\textbf{d)}
		
		Ao perder os dois eletrões \(3s\):
		
		\[
		\ce{Mg -> Mg^{2+} + 2e-}
		\]
		
		o magnésio adquire:
		
		\[
		1s^2\,2s^2\,2p^6.
		\]
		
		Esta configuração corresponde a uma camada de valência completa.
		
		A remoção de um terceiro eletrão exigiria muito mais energia, pois teria
		de ser removido um eletrão da camada \(n=2\).
		
	\end{resolucao}
	
	% ============================================================
	% DISTRIBUIÇÃO ELETRÓNICA
	% ============================================================
	
	\section*{\textcolor{eublue}{6. Distribuição eletrónica}}
	
	\begin{resolucao}{Escrever distribuições eletrónicas}
		
		\begin{center}
			
			\begin{tabular}{ccl}
				\toprule
				\textbf{Elemento}
				&
				\(\boldsymbol{Z}\)
				&
				\textbf{Distribuição eletrónica}
				\\
				\midrule
				
				\ce{C}
				&
				6
				&
				\(1s^2\,2s^2\,2p^2\)
				\\[2mm]
				
				\ce{O}
				&
				8
				&
				\(1s^2\,2s^2\,2p^4\)
				\\[2mm]
				
				\ce{Ne}
				&
				10
				&
				\(1s^2\,2s^2\,2p^6\)
				\\[2mm]
				
				\ce{Na}
				&
				11
				&
				\(1s^2\,2s^2\,2p^6\,3s^1\)
				\\[2mm]
				
				\ce{Mg}
				&
				12
				&
				\(1s^2\,2s^2\,2p^6\,3s^2\)
				\\[2mm]
				
				\ce{Al}
				&
				13
				&
				\(1s^2\,2s^2\,2p^6\,3s^2\,3p^1\)
				\\[2mm]
				
				\ce{Cl}
				&
				17
				&
				\(1s^2\,2s^2\,2p^6\,3s^2\,3p^5\)
				\\[2mm]
				
				\ce{Ca}
				&
				20
				&
				\(1s^2\,2s^2\,2p^6\,3s^2\,3p^6\,4s^2\)
				\\
				
				\bottomrule
			\end{tabular}
			
		\end{center}
		
	\end{resolucao}
	
	% ============================================================
	
	\begin{resolucao}{Diagramas de orbitais}
		
		\textbf{a) Carbono}
		
		\[
		1s^2\,2s^2\,2p^2
		\]
		
		Segundo a regra de Hund, os dois eletrões \(2p\) ocupam inicialmente
		orbitais diferentes.
		
		\begin{center}
			
			\begin{tikzpicture}[scale=0.9]
				
				\node at (-0.6,0.5) {\(1s\)};
				\draw (0,0) rectangle (1,1);
				
				\draw[-{Latex},thick,eublue]
				(0.30,0.15)--(0.30,0.85);
				
				\draw[{Latex}-,thick,eured]
				(0.70,0.15)--(0.70,0.85);
				
				\node at (1.8,0.5) {\(2s\)};
				\draw (2.4,0) rectangle (3.4,1);
				
				\draw[-{Latex},thick,eublue]
				(2.70,0.15)--(2.70,0.85);
				
				\draw[{Latex}-,thick,eured]
				(3.10,0.15)--(3.10,0.85);
				
				\node at (4.2,0.5) {\(2p\)};
				
				\foreach \x in {4.8,5.8,6.8}{
					\draw (\x,0) rectangle ({\x+1},1);
				}
				
				\draw[-{Latex},thick,eublue]
				(5.10,0.15)--(5.10,0.85);
				
				\draw[-{Latex},thick,eublue]
				(6.10,0.15)--(6.10,0.85);
				
			\end{tikzpicture}
			
		\end{center}
		
		Assim:
		
		\[
		\boxed{
			2p:
			\quad
			[\uparrow]\,
			[\uparrow]\,
			[\ ]
		}
		\]
		
		\vspace{4mm}
		
		\textbf{b) Oxigénio}
		
		\[
		1s^2\,2s^2\,2p^4
		\]
		
		Os três primeiros eletrões \(2p\) ocupam os três orbitais separadamente.
		O quarto eletrão emparelha num deles.
		
		\begin{center}
			
			\begin{tikzpicture}[scale=0.9]
				
				\node at (-0.6,0.5) {\(1s\)};
				\draw (0,0) rectangle (1,1);
				
				\draw[-{Latex},thick,eublue]
				(0.30,0.15)--(0.30,0.85);
				
				\draw[{Latex}-,thick,eured]
				(0.70,0.15)--(0.70,0.85);
				
				\node at (1.8,0.5) {\(2s\)};
				\draw (2.4,0) rectangle (3.4,1);
				
				\draw[-{Latex},thick,eublue]
				(2.70,0.15)--(2.70,0.85);
				
				\draw[{Latex}-,thick,eured]
				(3.10,0.15)--(3.10,0.85);
				
				\node at (4.2,0.5) {\(2p\)};
				
				\foreach \x in {4.8,5.8,6.8}{
					\draw (\x,0) rectangle ({\x+1},1);
				}
				
				% primeiro orbital: par
				\draw[-{Latex},thick,eublue]
				(5.08,0.15)--(5.08,0.85);
				
				\draw[{Latex}-,thick,eured]
				(5.48,0.15)--(5.48,0.85);
				
				% restantes
				\draw[-{Latex},thick,eublue]
				(6.10,0.15)--(6.10,0.85);
				
				\draw[-{Latex},thick,eublue]
				(7.10,0.15)--(7.10,0.85);
				
			\end{tikzpicture}
			
		\end{center}
		
		Assim:
		
		\[
		\boxed{
			2p:
			\quad
			[\uparrow\downarrow]\,
			[\uparrow]\,
			[\uparrow]
		}
		\]
		
	\end{resolucao}
	
	% ============================================================
	
	\begin{resolucao}{Detetar erros}
		
		Foram apresentadas:
		
		\[
		A:\quad1s^2\,2s^2\,2p^7
		\]
		
		\[
		B:\quad1s^2\,2s^2\,2p^6\,3s^3
		\]
		
		\[
		C:\quad1s^2\,2s^2\,2p^6\,3s^2\,3p^5
		\]
		
		\textbf{a)}
		
		São impossíveis:
		
		\[
		\boxed{A\text{ e }B}
		\]
		
		\textbf{b)}
		
		Um subnível \(p\) contém três orbitais e pode possuir, no máximo:
		
		\[
		3\times2=6
		\]
		
		eletrões.
		
		Assim:
		
		\[
		2p^7
		\]
		
		é impossível.
		
		Um subnível \(s\) possui apenas um orbital e pode conter, no máximo:
		
		\[
		2
		\]
		
		eletrões.
		
		Assim:
		
		\[
		3s^3
		\]
		
		é impossível.
		
		\textbf{c)}
		
		Na distribuição C:
		
		\[
		2+2+6+2+5=17
		\]
		
		Portanto:
		
		\[
		Z=17
		\]
		
		e o elemento é:
		
		\[
		\boxed{\ce{Cl}}
		\]
		
		\textbf{d)}
		
		A camada de valência é:
		
		\[
		3s^2\,3p^5
		\]
		
		Logo:
		
		\[
		2+5=7
		\]
		
		\[
		\boxed{7\text{ eletrões de valência}}
		\]
		
	\end{resolucao}
	
	% ============================================================
	
	\begin{resolucao}{Distribuição eletrónica de iões}
		
		\textbf{a) \(\ce{Na+}\)}
		
		O sódio neutro possui 11 eletrões:
		
		\[
		\ce{Na}:
		1s^2\,2s^2\,2p^6\,3s^1
		\]
		
		Ao perder um eletrão:
		
		\[
		\boxed{
			\ce{Na+}:
			1s^2\,2s^2\,2p^6
		}
		\]
		
		Possui 10 eletrões.
		
		\vspace{3mm}
		
		\textbf{b) \(\ce{Mg^{2+}}\)}
		
		O magnésio possui:
		
		\[
		1s^2\,2s^2\,2p^6\,3s^2
		\]
		
		Perdendo dois eletrões:
		
		\[
		\boxed{
			\ce{Mg^{2+}}:
			1s^2\,2s^2\,2p^6
		}
		\]
		
		Possui 10 eletrões.
		
		\vspace{3mm}
		
		\textbf{c) \(\ce{Cl-}\)}
		
		O cloro neutro possui:
		
		\[
		1s^2\,2s^2\,2p^6\,3s^2\,3p^5
		\]
		
		Ao receber um eletrão:
		
		\[
		\boxed{
			\ce{Cl-}:
			1s^2\,2s^2\,2p^6\,3s^2\,3p^6
		}
		\]
		
		Possui 18 eletrões.
		
		\vspace{3mm}
		
		\textbf{d) \(\ce{O^{2-}}\)}
		
		O oxigénio neutro possui:
		
		\[
		1s^2\,2s^2\,2p^4
		\]
		
		Ao receber dois eletrões:
		
		\[
		\boxed{
			\ce{O^{2-}}:
			1s^2\,2s^2\,2p^6
		}
		\]
		
		Possui 10 eletrões.
		
		\vspace{3mm}
		
		\textbf{e)}
		
		Na formulação original da ficha, \textbf{as quatro espécies não possuem
			a mesma distribuição eletrónica}.
		
		As espécies:
		
		\[
		\boxed{
			\ce{Na+},\quad
			\ce{Mg^{2+}},\quad
			\ce{O^{2-}}
		}
		\]
		
		possuem 10 eletrões e são isoeletrónicas com o néon:
		
		\[
		1s^2\,2s^2\,2p^6.
		\]
		
		Por outro lado:
		
		\[
		\boxed{\ce{Cl-}}
		\]
		
		possui 18 eletrões e é isoeletrónico com o árgon:
		
		\[
		1s^2\,2s^2\,2p^6\,3s^2\,3p^6.
		\]
		
		\begin{nota}
			
			Se se pretender que todas as espécies do exercício sejam isoeletrónicas
			com o néon, pode substituir-se \(\ce{Cl-}\), por exemplo, por
			\(\ce{F-}\).
			
			Nesse caso:
			
			\[
			\ce{F-}:
			1s^2\,2s^2\,2p^6
			\]
			
			e todas as espécies apresentariam 10 eletrões.
			
		\end{nota}
		
	\end{resolucao}
	
	% ============================================================
	% PROBLEMA INTEGRADOR
	% ============================================================
	
	\newpage
	
	\section*{\textcolor{euorange}{7. Problema integrador}}
	
	\begin{tcolorbox}[
		enhanced,
		breakable,
		colback=eulightorange,
		colframe=euorange,
		title=\textbf{Resolução},
		fonttitle=\bfseries
		]
		
		Dados:
		
		\[
		\lambda=\SI{486}{nm}
		=
		\SI{4,86e-7}{\meter}
		\]
		
		\[
		m(\ce{H})=\SI{2.00}{mg}
		=
		\SI{2.00e-3}{g}
		\]
		
		\[
		M(\ce{H})=\SI{1.008}{g\per mol}
		\]
		
		\textbf{a) Frequência}
		
		\[
		f=\frac{c}{\lambda}
		\]
		
		\[
		f=
		\frac{3,00\times10^8}
		{4,86\times10^{-7}}
		\]
		
		\[
		\boxed{
			f\approx\SI{6,17e14}{\hertz}
		}
		\]
		
		\vspace{3mm}
		
		\textbf{b) Energia de um fotão}
		
		\[
		E=hf
		\]
		
		\[
		E=
		6,626\times10^{-34}
		\times
		6,17\times10^{14}
		\]
		
		\[
		\boxed{
			E\approx\SI{4,09e-19}{\joule}
		}
		\]
		
		\vspace{3mm}
		
		\textbf{c) Energia no nível \(n=4\)}
		
		\[
		E_4=
		-\frac{2,18\times10^{-18}}{4^2}
		\]
		
		\[
		\boxed{
			E_4\approx\SI{-1,36e-19}{\joule}
		}
		\]
		
		\vspace{3mm}
		
		\textbf{d) Energia no nível \(n=2\)}
		
		\[
		E_2=
		-\frac{2,18\times10^{-18}}{2^2}
		\]
		
		\[
		\boxed{
			E_2=\SI{-5,45e-19}{\joule}
		}
		\]
		
		\vspace{3mm}
		
		\textbf{e) Variação de energia}
		
		\[
		\Delta E=E_2-E_4
		\]
		
		\[
		\Delta E
		=
		-5,45\times10^{-19}
		-
		(-1,36\times10^{-19})
		\]
		
		\[
		\boxed{
			\Delta E\approx\SI{-4,09e-19}{\joule}
		}
		\]
		
		\vspace{3mm}
		
		\textbf{f)}
		
		O módulo da variação da energia é:
		
		\[
		|\Delta E|
		=
		\SI{4,09e-19}{\joule}
		\]
		
		e a energia do fotão calculada anteriormente é:
		
		\[
		E_{\text{fotão}}
		=
		\SI{4,09e-19}{\joule}.
		\]
		
		Logo:
		
		\[
		\boxed{
			E_{\text{fotão}}
			=
			|\Delta E|
		}
		\]
		
		dentro do arredondamento utilizado.
		
		\vspace{3mm}
		
		\textbf{g)}
		
		Na transição:
		
		\[
		n=4\rightarrow n=2
		\]
		
		o átomo perde energia.
		
		Essa energia é transferida para o fotão emitido.
		
		Por conservação da energia:
		
		\[
		\boxed{
			E_{\text{fotão}}
			=
			|\Delta E_{\text{átomo}}|
		}
		\]
		
		\vspace{3mm}
		
		\textbf{h) Quantidade de matéria}
		
		\[
		n=\frac{m}{M}
		\]
		
		\[
		n=
		\frac{2,00\times10^{-3}}
		{1,008}
		\]
		
		\[
		\boxed{
			n\approx\SI{1,98e-3}{\mole}
		}
		\]
		
		\vspace{3mm}
		
		\textbf{i) Número de átomos}
		
		\[
		N=nN_A
		\]
		
		\[
		N=
		1,98\times10^{-3}
		\times
		6,022\times10^{23}
		\]
		
		\[
		\boxed{
			N\approx1,19\times10^{21}
			\text{ átomos}
		}
		\]
		
		\vspace{3mm}
		
		\textbf{j) Energia total emitida}
		
		No modelo idealizado, cada átomo emite um fotão.
		
		Assim:
		
		\[
		E_{\text{total}}
		=
		NE_{\text{fotão}}
		\]
		
		\[
		E_{\text{total}}
		=
		1,19\times10^{21}
		\times
		4,09\times10^{-19}
		\]
		
		\[
		E_{\text{total}}
		\approx
		\SI{4,89e2}{\joule}
		\]
		
		\[
		\boxed{
			E_{\text{total}}
			\approx
			\SI{489}{J}
		}
		\]
		
		Este valor corresponde ao modelo idealizado proposto no enunciado,
		no qual todos os átomos realizariam exatamente uma transição
		\(n=4\rightarrow n=2\).
		
	\end{tcolorbox}
	
	% ============================================================
	% ASTRONOMIA
	% ============================================================
	
	\section*{\textcolor{eugreen}{8. Ciência no mundo real}}
	
	\begin{resolucao}{Interpretar o espetro de uma estrela}
		
		No espetro da estrela foram identificadas linhas em:
		
		\[
		410\,\mathrm{nm},
		\quad
		434\,\mathrm{nm},
		\quad
		486\,\mathrm{nm},
		\quad
		656\,\mathrm{nm}.
		\]
		
		\textbf{a)}
		
		Estes valores coincidem aproximadamente com linhas características do
		espetro do átomo de hidrogénio na série de Balmer.
		
		\textbf{b)}
		
		Pode concluir-se que:
		
		\[
		\boxed{
			\text{existe hidrogénio na região da estrela responsável pela absorção}
		}
		\]
		
		\textbf{c)}
		
		Os átomos de hidrogénio absorvem determinados comprimentos de onda,
		associados às diferenças entre os seus níveis eletrónicos de energia.
		
		Essas linhas ficam registadas na luz recebida da estrela.
		
		Assim, analisando a luz com um espetroscópio e comparando o padrão
		observado com espetros obtidos em laboratório, é possível identificar
		o hidrogénio sem recolher diretamente uma amostra da estrela.
		
	\end{resolucao}
	
	% ============================================================
	% SÍNTESE
	% ============================================================
	
	\section*{\textcolor{eublue}{9. Síntese das ideias fundamentais}}
	
	\begin{tcolorbox}[
		colback=eugray,
		colframe=eudark,
		breakable
		]
		
		\textbf{Quantidade de matéria}
		
		\[
		n=\frac{m}{M}
		\qquad
		N=nN_A
		\]
		
		\vspace{3mm}
		
		\textbf{Energia da radiação}
		
		\[
		E=hf
		\]
		
		\[
		c=\lambda f
		\]
		
		\[
		E=\frac{hc}{\lambda}
		\]
		
		\vspace{3mm}
		
		\textbf{Átomo de hidrogénio}
		
		\[
		E_n=
		-\frac{2,18\times10^{-18}}{n^2}
		\ \mathrm{J}
		\]
		
		\[
		\Delta E=E_f-E_i
		\]
		
		Se:
		
		\[
		\Delta E<0
		\]
		
		ocorre emissão.
		
		Se:
		
		\[
		\Delta E>0
		\]
		
		ocorre absorção.
		
		Em qualquer transição radiativa:
		
		\[
		E_{\text{fotão}}
		=
		|\Delta E|
		\]
		
		\vspace{3mm}
		
		\textbf{Espetros}
		
		Cada elemento possui um conjunto característico de níveis eletrónicos
		e, consequentemente, um espetro característico.
		
		\vspace{3mm}
		
		\textbf{Espetro fotoeletrónico}
		
		A energia de remoção permite obter informação sobre a forma como os
		eletrões estão distribuídos pelos diferentes níveis e subníveis.
		
		\[
		E_{\text{fotão}}
		=
		E_{\text{remoção}}
		+
		E_c
		\]
		
		\vspace{3mm}
		
		\textbf{Distribuição eletrónica}
		
		O preenchimento dos orbitais obedece ao:
		
		\begin{itemize}
			\item princípio da energia mínima;
			\item princípio da exclusão de Pauli;
			\item regra de Hund.
		\end{itemize}
		
	\end{tcolorbox}
	
	\vfill
	
	% ============================================================
	% RODAPÉ VISUAL
	% ============================================================
	
	\begin{center}
		
		\begin{tikzpicture}
			
			\draw[eublue,very thick]
			(0,0)--(13,0);
			
			\node[
			text=eublue,
			font=\Large\bfseries
			]
			at (6.5,0.55)
			{EUFISICA exercícios};
			
			\node[
			text=eudark,
			font=\small
			]
			at (6.5,-0.45)
			{Resolução -- Física e Química A -- 10.º ano};
			
			\node[
			text=eudark,
			font=\small
			]
			at (6.5,-0.85)
			{Prof. José Gonçalves};
			
		\end{tikzpicture}
		
	\end{center}
	
\end{document}